\documentclass[10pt]{article}

%% Esempio 4 - Verifiche scritte con fancylessons.
%% Con l'opzione legacy i vecchi comandi funzionano senza modifiche.
%% Per i testi nuovi conviene la sintassi degli esempi 1-3.
\usepackage[legacy]{fancytests}
\usepackage{lmodern}
\renewcommand*\familydefault{\sfdefault}
\usepackage[T1]{fontenc}

\geometry{a4paper, top=15mm, bottom=15mm, left=17.5mm, right=17.5mm}
\linespread{1.25}

\begin{document}

\mathtestheader{4M}{20/03/2026}{Logaritmo, equazioni esponenziali e logaritmiche}()

\exercise Risolvi le seguenti equazioni.

\begin{splitlist}(1.5)%
    \itemx $3^x - 2^x = 0$
        & \score{1.0} &
    \itemx $\log_{10}(2x - 1) + \log_{10}(x) = 1$
        & \score{1.0}\\
    \itemx $3 \log_5(x - 2) = -3$
        & \score{1.0} &
    \itemx $e^{2x} - e^{x + 1} = 0$
        & \score{1.0}
\end{splitlist}

\exercisescore

\exercise Rispondi ai seguenti quesiti.

\begin{scoredlist}(1.5)<L>%
    \itemx Rispetto alla funzione $f(x) = \ln(x)$, come risulta il grafico di
        \randOp{$g(x) = \ln(x + 3)$}{$g(x) = \ln(x - 3)$}?
        & \score{1.0}\\[-1em]
        & \randomfouroptions{%
            Traslato verso l'alto di 3 unità.
        }{%
            Traslato verso destra di 3 unità.
        }{%
            Traslato verso il basso di 3 unità.
        }{%
            Traslato verso sinistra di 3 unità.
        }(.5)<L>|0pt|\\[1em]
    \itemx Quale delle seguenti funzioni logaritmiche è rappresentata?
        & \score{1.0}\\[-1em]
        & \randomthreeplusoneoptions{%
            $f(x) = \log_3(x)$
        }{%
            $f(x) = \log_2(x)$
        }{%
            $f(x) = \ln(x)$
        }{%
            Nessuna
        }(.5)<L>|0pt|\\
\end{scoredlist}

\exercisescore

\exercise Spiega la differenza tra scala lineare e scala logaritmica.

\exercisescore[2.0]

\vfill
\minimalfinalscore

\vspace{-2.25em}
\finalscorenotable

\end{document}
